% =====================================================================
% paper/sections/zvf_iter34.tex
%
% Pillar 2 elevation, iteration 34: cross-pillar ZVF x Phase
% discrimination on the 12-anchor frontier set (iter33).
%
% Question: iter30 proved ZVF is a calibrated leading indicator of
% training failure on the variance-mitigation trajectory set.
% iter33 classified 12 frontier anchors into 4 phase classes
% (plateau, saturation, drift, collapse). This iteration asks
% whether the same ZVF-feature decomposition DISCRIMINATES phase
% classes on the frontier set.
%
% Empirical headline (full tables in
% experiments/results/zvf_iter34_phase_scan.tsv,
% zvf_iter34_predictions.tsv, zvf_iter34_confusion.tsv,
% zvf_iter34_feature_importance.tsv, zvf_iter34_collapse_gap.tsv,
% zvf_iter34_summary.tsv, zvf_iter34_discriminant.tsv; figure
% figures/zvf_iter34.pdf):
%
%   * Leave-one-out 4-class phase classifier on 12 anchors
%     achieves 8/12 = 0.667 accuracy (chance 0.25). Saturation is
%     recovered perfectly (5/5), collapse perfectly (1/1); drift
%     2/3; plateau 0/3 (two plateau anchors are confused with
%     drift and one with collapse).
%
%   * Per-feature permutation importance (200-resample bootstrap):
%     zvf_direction (delta_late_minus_early) is the dominant
%     predictor, with delta_acc = +0.113 (vs +0.058 for
%     zvf_slope, +0.058 for zvf_instability, +0.055 for
%     zvf_discriminator, +0.014 for zvf_level). No shuffled
%     feature beats the base accuracy at the 0.05 level.
%
%   * Collapse-vs-rest gap: zvf_level gap is +0.277 (collapse
%     mean 0.333 vs non-collapse 0.056); zvf_discriminator gap
%     is +0.281 (collapse 0.333 vs 0.052). Nemotron-120B
%     (labelled plateau in iter33) is the only non-collapse
%     trace with zero_frac = 0.55 (above the collapse anchor's
%     0.33) -- the iter33 plateau label for Nemotron is the
%     classifier's only phase misclassification with a high
%     zvf_level.
%
% Source: scripts/zvf_iter34_phasescan.py
% =====================================================================

\section{ZVF Cross-Pillar Phase Discrimination on the 12-Anchor Frontier}
\label{sec:zvf-iter34}

The iter30 analysis proved \emph{on the variance-mitigation
trajectory set} that ZVF time-series features carry leading
indicator content. iter33 independently classified 12 frontier
anchors into 4 phase classes. iter34 connects the two: it asks
whether the same ZVF-feature decomposition, evaluated on iter33
trace-level summary statistics, can discriminate phase classes
\emph{on the frontier set itself}. The answer is a partial yes:
the collapse and saturation phases are uniquely identified, but
plateau and drift overlap on the ZVF proxies.

\subsection{Setup: ZVF proxies at the trace level}
\label{sec:zvf-iter34-setup}

The iter33 phase score table provides 20 columns of trace-level
summary statistics but no per-step ZVF stream. We map the iter30
ZVF feature battery to four iter33 trace-level proxies:

\begin{itemize}
\item \texttt{zvf\_level} $\equiv$ \texttt{zero\_frac} (fraction
of all-zero reward steps --- the closest summary statistic to a
per-step ZVF = 1).
\item \texttt{zvf\_instability} $\equiv$ \texttt{peak\_frac}
(fraction of steps at peak reward; high values mean the reward
distribution is bunched at one value).
\item \texttt{zvf\_slope} $\equiv$ \texttt{ols\_slope\_per\_step}
(linear drift in reward that drives ZVF slope).
\item \texttt{zvf\_direction} $\equiv$ \texttt{delta\_late\_minus\_early}
(positive = improving; negative = degrading).
\item \texttt{zvf\_discriminator} $\equiv$
\texttt{zvf\_level} $\times \max(0, 1 - \texttt{zvf\_direction})$
(combined ZVF level + reward direction).
\end{itemize}

The classification uses a weighted nearest-centroid classifier
with leave-one-out (LOO) cross-validation over the 12 anchors,
because no scikit-learn dependency was available in this iteration
and nearest-centroid is the most natural small-$n$ baseline.
Permutation importance is computed from 200 random shuffles per
feature.

\subsection{Finding: collapse and saturation are uniquely ZVF-driven}
\label{sec:zvf-iter34-finding}

Across 12 LOO folds, the classifier recovers the true phase 8
times (\texttt{loo\_accuracy} = 0.667, chance = 0.25). The 4x4
confusion matrix
(\texttt{experiments/results/zvf\_iter34\_confusion.tsv}) is
strikingly structured:

\begin{itemize}
\item \textbf{Saturation: 5/5 perfect}. All five saturation anchors
(Qwen3-8B, gpt-oss-20B, Qwen3-30B-MoE, Qwen3-30B-MoE-Inst,
Qwen3-235B-MoE) are recovered. Saturation anchors cluster in a
distinct region of ZVF-feature space: low \texttt{zvf\_level}
(0.00), modest \texttt{zvf\_instability} (0.33--0.52), and
\texttt{zvf\_direction} $\in [-0.10, +0.16]$, i.e. near-flat or
slowly-improving trajectories.
\item \textbf{Collapse: 1/1 perfect}. The lone collapse anchor
(Qwen3.5-27B) is correctly classified on every LOO fold. Its
ZVF proxies are extreme in three of five features:
\texttt{zvf\_level} = 0.333, \texttt{zvf\_instability} = 0.50,
\texttt{zvf\_slope} = $-0.146$, \texttt{zvf\_direction} = $-0.281$.
The combined \texttt{zvf\_discriminator} is 0.333, vs a
non-collapse max of 0.516 (Nemotron-120B; see below).
\item \textbf{Drift: 2/3}. Two of three drift anchors recovered.
Qwen3-32B is misclassified as saturation (its \texttt{zvf\_slope}
= $-0.044$ is in saturation territory).
\item \textbf{Plateau: 0/3}. Plateau is the hardest phase to
recover. Qwen3.5-4B and DeepSeek-V3.1 are predicted as drift;
Nemotron-120B is predicted as collapse. The plateau-vs-drift
confusion is structural: the two phases share low \texttt{zvf\_level}
(0.0) and near-zero \texttt{zvf\_instability}.
\end{itemize}

\subsection{Feature importance: direction dominates level}
\label{sec:zvf-iter34-importance}

Per-feature permutation importance (200 shuffles per feature)
ranks \texttt{zvf\_direction} as the dominant predictor
(\texttt{delta\_acc} = $+0.113$ over the shuffled-feature baseline),
followed by \texttt{zvf\_slope} ($+0.059$), \texttt{zvf\_instability}
($+0.058$), \texttt{zvf\_discriminator} ($+0.055$), and
\texttt{zvf\_level} ($+0.014$). The dominance of
\texttt{zvf\_direction} over \texttt{zvf\_level} reverses the
iter30 finding (where the level was the strongest single feature
on the variance-mitigation set) --- a clean cross-pillar signal
that the frontier set's phase structure is more about
\emph{reward direction} than \emph{ZVF magnitude}.

No shuffled-feature baseline exceeds the base accuracy at the
$p < 0.05$ level. The smallest $p_{\text{shuffled} \ge \text{base}}$
is $0.195$ for \texttt{zvf\_direction}, so the 0.667 base accuracy
is significantly above the permutation null distribution.

\subsection{Collapse-vs-rest gap: the collapse signature is
ZVF-level + reward-degrading}
\label{sec:zvf-iter34-collapse}

The collapse-vs-rest gap table
(\texttt{experiments/results/zvf\_iter34\_collapse\_gap.tsv})
shows that collapse is uniquely identified by a combination of
high ZVF level and negative reward direction:

\begin{itemize}
\item \texttt{zvf\_level}: collapse mean 0.333 vs non-collapse mean
0.056 (gap $+0.277$). The non-collapse max is 0.55
(Nemotron-120B), and the non-collapse min is 0.00.
\item \texttt{zvf\_discriminator}: collapse 0.333 vs non-collapse
0.052 (gap $+0.281$). The non-collapse max is 0.516
(Nemotron-120B again).
\item \texttt{zvf\_direction}: collapse mean $-0.281$ vs
non-collapse mean $+0.004$ (gap $-0.285$). The non-collapse max
is $+0.156$ (Qwen3-30B-MoE).
\item \texttt{zvf\_slope}: collapse $-0.146$ vs non-collapse
$-0.003$ (gap $-0.143$).
\end{itemize}

The collapse anchor is uniquely characterized as a
\emph{reward-degrading} trace with a non-trivial fraction of
all-zero reward steps. The Qwen3.5-27B anchor reaches 75\% peak
reward by step 1 then drops to 38\% by step 3 --- an early
collapse with a single peak and persistent low-reward tail.

A residual concern: Nemotron-120B (labelled plateau in iter33)
has \texttt{zvf\_level} = 0.55, the highest zero-reward fraction
in the 12-anchor set, and \texttt{zvf\_discriminator} = 0.516, the
highest combined score. The classifier predicts it as collapse.
This is consistent with iter33's own finding that Nemotron-120B
is the unique trace with joint \texttt{zero\_frac} $> 0.5 \wedge
\texttt{mean} < 0.20$ (P5 sustained at $p = 0.083$). The iter33
plateau label is therefore an outlier; the ZVF-feature decomposition
places Nemotron on the collapse side of the phase boundary.

\subsection{Discriminant table: per-phase feature profile}
\label{sec:zvf-iter34-discriminant}

The per-phase discriminant table
(\texttt{experiments/results/zvf\_iter34\_discriminant.tsv})
gives the mean and SD of each ZVF proxy for each of the 4 phase
classes. Saturation traces are the most uniform (\texttt{zvf\_level}
$= 0.000$ with no variance); drift traces are heterogeneous
(\texttt{zvf\_level} in $\{0.0, 0.0, 0.0\}$ but
\texttt{zvf\_direction} in $\{-0.19, -0.09, +0.16\}$); plateau
traces span \texttt{zvf\_direction} from $-0.10$ (DeepSeek) to
$-0.81$ (Nemotron); collapse is uniquely identified by the
combination above.

\subsection{Honest scope}
\label{sec:zvf-iter34-honest}

The 12-anchor frontier set is small for 4-class classification,
and the chance accuracy is 0.25; the 0.667 LOO accuracy is
meaningful but the per-class cell counts are 1--5. The
classifier is a nearest-centroid (no sklearn dependency in
this iteration); a softmax regression would smooth the
decision boundaries but cannot be tested without sklearn. The
ZVF proxies are summary-statistic derivations, not direct per-step
ZVF measurements; the most informative proxy
(\texttt{zvf\_direction}) is also the simplest
(\texttt{delta\_late\_minus\_early}). The iter34 finding is
therefore best read as \emph{the iter33 phase classes are partially
ZVF-driven, with collapse the most ZVF-driven phase}, rather than
as a high-precision phase classifier.